\chapter{Burnout}

\section*{Definition and Overview}
Burnout is a state of physical, emotional, and mental exhaustion caused by prolonged or excessive stress, usually related to work, although it can also arise from demanding caregiving, study, or other sustained responsibilities. It was first described in the caring professions and is now recognised by the World Health Organization as an occupational phenomenon, though not as a medical condition in itself.

People experiencing burnout typically feel drained and used up, increasingly detached or cynical about their work, and unable to cope or perform effectively. Burnout is distinct from depression but closely related to it, causes genuine distress and functional impairment, and should always be taken seriously. It is important to distinguish workplace burnout from clinical conditions such as depression and anxiety so that the most appropriate help is provided.

\section*{Epidemiology}
Burnout is common. Surveys suggest that a substantial proportion of working adults experience significant burnout at some stage of their careers, with higher rates in professions involving intense human interaction, such as health care, teaching, social work, and emergency services. It is more frequent in younger and mid-career workers, in people with heavy workloads and low control over how work is done, and in those who feel a strong sense of responsibility for others' welfare. Women report burnout somewhat more often than men, partly reflecting the combined load of paid work and unpaid caregiving. Reported rates rose during and after the COVID-19 pandemic, particularly among healthcare workers, and burnout remains a leading reason for sickness absence and early career change.

\section*{Aetiology and Pathophysiology}
Burnout arises from a sustained imbalance between the demands placed on a person and the resources available to meet them. Contributing factors include:
\begin{itemize}
    \item \textbf{Chronic workload:} excessive hours, unrelenting deadlines, and understaffing
    \item \textbf{Lack of control:} little say over how, when, and where work is performed
    \item \textbf{Insufficient reward:} inadequate pay, recognition, or satisfaction relative to effort
    \item \textbf{Unfairness and conflict:} perceived unfair treatment, workplace bullying, or clashes in values
    \item \textbf{Poor social support:} isolation, unsupportive management, or weak team relationships
    \item \textbf{Personal factors:} perfectionism, difficulty setting boundaries, and over-identification with work
    \item \textbf{Lack of recovery:} long commutes, shift work, and inadequate rest, leave, and downtime
\end{itemize}

The underlying mechanism is chronic stress. Prolonged activation of the body's stress response, with sustained elevation of stress hormones, is thought to exhaust energy reserves, disrupt sleep, mood, and immune function, and alter the way the brain processes reward and motivation. Over time this produces the classic triad of exhaustion, cynicism, and reduced personal efficacy.

\section*{Clinical Features}
\begin{itemize}
    \item \textbf{Exhaustion:} deep fatigue that rest and weekends do not relieve
    \item \textbf{Cynicism and detachment:} feeling negative, callous, or disconnected from work and colleagues
    \item \textbf{Reduced effectiveness:} difficulty concentrating, forgetfulness, and feeling that achievements are worthless
    \item \textbf{Sleep problems:} difficulty falling or staying asleep, or sleeping excessively
    \item \textbf{Physical symptoms:} headaches, muscle tension, digestive disturbance, and frequent minor illnesses
    \item \textbf{Emotional changes:} irritability, low mood, anxiety, and loss of motivation or enjoyment
    \item \textbf{Behavioural changes:} withdrawal from colleagues, avoidance of work, increased use of alcohol or comfort food, and reduced self-care
\end{itemize}

\section*{Differential Diagnosis}
\begin{itemize}
    \item \textbf{Major depressive disorder:} overlapping symptoms such as low mood, sleep disturbance, and loss of interest, but a distinct, more pervasive condition requiring specific treatment
    \item \textbf{Anxiety disorders:} generalised anxiety, panic, or social anxiety that may coexist with burnout
    \item \textbf{Adjustment disorder:} distress arising in response to a specific identifiable stressful event
    \item \textbf{Chronic fatigue syndrome:} where fatigue is the dominant symptom and is typically worsened by exertion
    \item \textbf{Medical causes of fatigue:} thyroid disease, anaemia, vitamin deficiency, diabetes, or sleep apnoea
    \item \textbf{Substance use:} alcohol or stimulant misuse that contributes to or masks exhaustion
\end{itemize}

\section*{When to Seek Medical Care}
People should see a doctor if exhaustion, low mood, or anxiety persists for more than a few weeks, is affecting sleep, health, or relationships, or is interfering with work or study. Early help can prevent burnout from progressing into depression or anxiety disorders.

\textbf{Seek immediate help if there are thoughts of suicide or self-harm.} Call 000 (or local emergency services), contact a crisis support line, or go to the nearest hospital emergency department without delay.

\section*{Diagnosis and Investigations}
\begin{itemize}
    \item \textbf{Clinical interview:} a full history of work demands, stress, sleep, mood, and day-to-day functioning; validated burnout questionnaires such as the Maslach Burnout Inventory may be used in research or workplace settings
    \item \textbf{Workplace assessment:} exploration of workload, control, support, and any specific workplace issues or conflicts
    \item \textbf{Screening for depression and anxiety:} standard questionnaires such as the PHQ-9 and GAD-7 help distinguish burnout from clinical disorders
    \item \textbf{Physical examination and blood tests:} to exclude medical causes of fatigue, including full blood count, thyroid function, iron studies, vitamin B12, and blood glucose
    \item \textbf{Sleep assessment:} referral for sleep studies if sleep apnoea or another sleep disorder is suspected
\end{itemize}

\section*{Management}
\begin{itemize}
    \item \textbf{Address the source:} reduce workload where possible, set boundaries, delegate tasks, and negotiate changes with employers
    \item \textbf{Recovery time:} take leave, schedule genuine rest, and protect time away from work and email
    \item \textbf{Physical health:} prioritise regular sleep, physical activity, healthy eating, and reduced alcohol and caffeine
    \item \textbf{Psychological support:} counselling and cognitive behaviour therapy (CBT) help manage stress and rebuild coping skills
    \item \textbf{Workplace measures:} discuss reasonable adjustments with managers or human resources; employee assistance programs offer confidential counselling
    \item \textbf{Treatment of coexisting problems:} depression or anxiety, if present, should be treated according to a doctor's advice
    \item \textbf{Follow-up:} regular review with a general practitioner to monitor progress and prevent relapse
\end{itemize}

\section*{Complications}
\begin{itemize}
    \item Depression and anxiety disorders, which are the most important downstream risks
    \item Sleep disorders, including chronic insomnia
    \item Physical health problems linked to chronic stress, including high blood pressure and cardiovascular disease
    \item Increased alcohol or drug use
    \item Relationship strain, social withdrawal, and isolation
    \item Sickness absence, job loss, and financial difficulty
    \item Reduced performance and a higher risk of errors at work, which is particularly important in safety-critical roles
\end{itemize}

\section*{Prognosis and Outcomes}
Burnout usually improves when the underlying causes are addressed and recovery time is taken. Many people return to full functioning within weeks to months, especially when they act early and make sustainable changes. If burnout persists, the risk of developing clinical depression increases and recovery becomes slower. Gradual return to work and lasting adjustments to workload and boundaries help prevent recurrence. People who seek help early and use both workplace and personal supports generally achieve the best outcomes.

\section*{Prevention and Self-Care}
\begin{itemize}
    \item Set realistic goals and learn to decline excessive demands
    \item Protect sleep, physical activity, and time for rest and hobbies
    \item Build and maintain supportive relationships at work and at home
    \item Take regular breaks during the day and use annual leave fully
    \item Develop stress-management skills such as mindfulness and relaxation techniques
    \item Reduce reliance on alcohol, caffeine, and other short-term crutches
    \item Discuss workload concerns with a manager early, and consult a doctor if symptoms persist
    \item Watch for early warning signs -- exhaustion, irritability, and withdrawal -- and act on them promptly
\end{itemize}

\section*{Sources and Further Reading}
The following sources provide reliable, accessible information on burnout and its management:
\begin{itemize}
    \item World Health Organization. \textit{Burn-out is an occupational phenomenon.} https://www.who.int/
    \item NHS. \textit{Stress and burnout.} https://www.nhs.uk/mental-health/
    \item Mayo Clinic. \textit{Burnout: job stress and your health.} https://www.mayoclinic.org/
    \item MedlinePlus. \textit{Burnout.} https://medlineplus.gov/
    \item Beyond Blue. \textit{Workplace mental health.} https://www.beyondblue.org.au/
    \item NICE. \textit{Workplace health: management of stress.} https://www.nice.org.uk/
\end{itemize}
